\section{Orbit algebras and the H\'enon composition}
\label{sec:orbit-algebras}

The quadratic H\'enon algebra basis, its finite and infinite versions, and the
wrapping description of the dynamics are due to
\citet{Bousch1992HenonAlgebras}.  The construction below extends that
normal-form argument to the prescribed phase degrees.  We give the reduction
and the identifications explicitly because both the ordinary-degree filtration
and the scheme-valued periodic test will use the actual basis, not merely a
generating family.  The orbit-coefficient obstruction obtained from this basis
in the next section is a linear-algebra consequence of its permutation action.

Throughout, $K$ is a field of characteristic zero and
$A=K[x,y]$.  Fix $k\ge1$ and polynomials $p_0,\ldots,p_{k-1}$ of degrees
$d_i\ge2$, and extend the subscripts periodically to $\Z$.  Write
\begin{equation}
  H_i(x,y)=(p_i(x)-y,x),\qquad
  F=H_{k-1}\circ\cdots\circ H_0,\qquad
  \sigma=F^*,\qquad \delta=\prod_{i=0}^{k-1}d_i.
  \label{eq:henon-composition}
\end{equation}
Here $F^*f=f\circ F$.  The inverse of $H_i$ is
$(x,y)\mapsto(y,p_i(y)-x)$, so these maps and $F$ are polynomial
automorphisms.  They preserve $\dd x\wedge\dd y$, since
$\dd(p_i(x)-y)\wedge\dd x=\dd x\wedge\dd y$.
No normalization of their nonzero leading coefficients is imposed.

The orbit coordinates are the polynomials determined by
\begin{equation}
  X_0=x,\qquad X_{-1}=y,\qquad
  p_i(X_i)=X_{i-1}+X_{i+1}\quad(i\in\Z).
  \label{eq:orbit-recurrence}
\end{equation}
The recurrence defines them in both directions: solve for $X_{i+1}$
starting at $i=0$, and for $X_{i-1}$ starting at $i=-1$.
An exponent word is a finitely supported sequence
$e=(e_i)_{i\in\Z}$ of integers with $0\le e_i<d_i$.
For such a word put
\begin{equation}
  M_e=\prod_{i\in\Z}X_i^{e_i},\qquad
  \cB=\{M_e:0\le e_i<d_i,\ e\text{ finitely supported}\}.
  \label{eq:standard-orbit-monomials}
\end{equation}
The zero word gives $M_0=1$.

\begin{theorem}[Orbit basis]
\label{thm:orbit-basis}
The family $\cB$ is a $K$-basis of $A$.  More precisely, the evaluation map
$Z_i\mapsto X_i$ identifies $A$ with
\[
  A_\infty=
  K[Z_i:i\in\Z]\big/
  (p_i(Z_i)-Z_{i-1}-Z_{i+1}:i\in\Z),
\]
and the residue classes of the monomials $\prod_iZ_i^{e_i}$ with
$0\le e_i<d_i$ form a basis of this quotient.
\end{theorem}

\begin{proof}
Write $p_i(T)=a_iT^{d_i}+q_i(T)$, where $a_i\in K^*$ and
$\deg q_i<d_i$.  In $R_\infty=K[Z_i:i\in\Z]$ consider the rules
\begin{equation}
  Z_i^{d_i}\longrightarrow r_i,
  \qquad
  r_i=a_i^{-1}\bigl(Z_{i-1}+Z_{i+1}-q_i(Z_i)\bigr).
  \label{eq:pure-power-rewrite}
\end{equation}
Every polynomial in $R_\infty$ involves only finitely many variables.
Although a rule can introduce a neighboring index that was absent from the
input, this does not prevent termination.  Replacing a factor $Z_i^{d_i}$
in a monomial of total degree $m$ produces only monomials of degree at most
$m-1$.  A recursively expanded monomial therefore has a finitely branching
reduction tree of depth at most $m$.  Arbitrary orders of polynomial
reduction terminate as well.  If the input has degree at most $D$, list the
numbers of its current nonzero monomial terms in degrees
$D,D-1,\ldots,0$, in this order.  Every reduction strictly decreases this
finite tuple in lexicographic order: one term at the degree being reduced
is removed, and only lower degrees can receive new terms.  Cancellations
can only decrease the tuple further.  Lexicographic descent on a fixed
finite product of copies of $\Z_{\ge0}$ is well founded.

We next prove uniqueness of reduction by induction on the degree of a
monomial.  Irreducible monomials require no choice.  If a monomial admits
two reductions by the same rule, the resulting polynomial is the same,
since selecting a different copy of the identical pure-power factor
does not change the quotient monomial.  For different rules $i\ne j$,
their pure powers are relatively prime.  Thus a monomial $M$ divisible
by both can be reduced, first at $i$ and then using the retained factor
$Z_j^{d_j}$ in every expanded term, to
\begin{equation}
  \frac{M}{Z_i^{d_i}Z_j^{d_j}}\,r_ir_j.
  \label{eq:common-multiple-reduction}
\end{equation}
Reversing the order gives the identical polynomial.  This remains true
when the indices are adjacent and $r_i$ contains $Z_j$: the original
$Z_j^{d_j}$ factor is still present before the second replacement.
All monomials after the first step have degree smaller than $\deg M$,
so the induction hypothesis makes their subsequent complete reductions
unique.  Consequently the two first choices have the same normal form.

Extend this normal form linearly and denote it by $\NF_\infty$.
The preceding argument also handles a reduction performed in any one
term of a polynomial.  For every monomial $M$ and every $i$ it gives
\[
  \NF_\infty\bigl(M(Z_i^{d_i}-r_i)\bigr)=0.
\]
Hence $\NF_\infty$ annihilates the relation ideal.  Conversely, a
reduction changes a polynomial by an element of that ideal and fixes
each irreducible monomial.  Every class is therefore spanned by the
bounded-exponent monomials, and a linear combination of them lying in
the ideal must be zero: applying $\NF_\infty$ leaves that combination
unchanged.  This proves linear independence as well as spanning.

It remains to identify the quotient with the specified phase space.
The recurrence \eqref{eq:orbit-recurrence} gives a homomorphism
\[
  \Psi:A_\infty\longrightarrow A,\qquad [Z_i]\longmapsto X_i.
\]
In the reverse direction define
\[
  \Theta:A\longrightarrow A_\infty,\qquad
  x\longmapsto[Z_0],\quad y\longmapsto[Z_{-1}].
\]
The composite $\Psi\Theta$ fixes $x$ and $y$.
Inside $A_\infty$, solving the defining relations successively to the
right and to the left expresses each $[Z_i]$ as
$X_i([Z_0],[Z_{-1}])$.  Therefore $\Theta\Psi$ fixes every generator
$[Z_i]$ and is the identity as well.  These homomorphisms are inverse,
and the abstract basis becomes precisely \eqref{eq:standard-orbit-monomials}.
\end{proof}

For $g\in A$, its normal form means its unique finite expansion in $\cB$.
We identify the abstract bounded-exponent monomials with their polynomial
images from now on.  In particular, the coefficient of $M_0=1$ refers to
this expansion and need not be the constant coefficient in the original
variables $x,y$.

\begin{proposition}[Macro shift]
\label{prop:macro-shift}
For every $i\in\Z$,
\begin{equation}
  \sigma X_i=X_{i+k}.
  \label{eq:macro-shift}
\end{equation}
Thus $\sigma$ permutes $\cB$ by shift-$k$, and normal-form reduction
commutes with this shift.
\end{proposition}

\begin{proof}
Successive application of $H_0,\ldots,H_{k-1}$ sends $(x,y)$ to
$(X_k,X_{k-1})$.  This proves \eqref{eq:macro-shift} for $i=0,-1$.
Applying $\sigma$ to \eqref{eq:orbit-recurrence} shows that
$(\sigma X_i)_{i\in\Z}$ satisfies the same recurrence as
$(X_{i+k})_{i\in\Z}$, because $p_{i+k}=p_i$.
The two adjacent initial values determine all other values in both
directions, which proves the formula for every $i$.
Since $d_{i+k}=d_i$, a shifted admissible word is still admissible.
The shift also preserves the relations in
\eqref{eq:pure-power-rewrite}; uniqueness of their normal forms proves
the last assertion.
\end{proof}

The shift in this proposition is attached to the chosen composition.
Even if some or all of the phase polynomials coincide, replacing
shift-$k$ by shift-$1$ changes the cohomological equation for $F$.

For periodic data we need a finite version of the same algebra.  Let
$n\ge1$ and put $N=kn\ge3$.  Since $k$ divides $N$, the phase data
are well defined on $\Z/N\Z$.  Define
\begin{equation}
  A_N=K[Z_i:i\bmod N]\big/
  (p_i(Z_i)-Z_{i-1}-Z_{i+1}:i\bmod N).
  \label{eq:finite-orbit-algebra}
\end{equation}
All subscripts in this presentation are cyclic.  Write
$\sigma_N[Z_i]=[Z_{i+k}]$ for the induced automorphism.

\begin{proposition}[Periodic algebra]
\label{prop:periodic-algebra}
For $N=kn\ge3$, the algebra $A_N$ has basis
\begin{equation}
  \cB_N=
  \left\{\prod_{i=0}^{N-1}[Z_i]^{e_i}:0\le e_i<d_i\right\}
  \quad\text{and}\quad
  \dim_K A_N=\delta^n.
  \label{eq:finite-basis-length}
\end{equation}
There is a surjective wrapping homomorphism $W_N:A\to A_N$ with
$W_N(X_i)=[Z_{i\bmod N}]$, and
\begin{equation}
  \ker W_N=J_N:=(X_N-X_0,\ X_{N-1}-X_{-1}).
  \label{eq:fixed-scheme-ideal}
\end{equation}
Consequently
\begin{equation}
  A_N\simeq A/J_N=K[\Fix(F^n)],\qquad
  W_N\sigma=\sigma_NW_N,\qquad \sigma_N^n=1.
  \label{eq:fixed-scheme-identification}
\end{equation}
These are identities for the complete fixed-point scheme; no reducedness
assumption is needed.
\end{proposition}

\begin{proof}
The cyclic version of \eqref{eq:pure-power-rewrite} still replaces
each pure power by terms of strictly smaller total degree.  Its distinct
leading pure powers remain relatively prime, and the common-multiple
calculation \eqref{eq:common-multiple-reduction} is unchanged.
The termination, uniqueness, and independence argument in the proof of
Theorem~\ref{thm:orbit-basis} therefore applies to this finite set of
variables.  It gives $\cB_N$ as a basis, with cardinality
$\prod_{i=0}^{N-1}d_i=\delta^n$.

Sending every infinite index to its residue modulo $N$ respects all
relations and defines $W_N$.  Each finite generator is in its image,
so it is surjective.  Both generators of $J_N$ map to zero.
To prove that they give the entire kernel, work first in $A/J_N$.
The sequences $(X_{i+N})$ and $(X_i)$ have the same two adjacent
values at $i=0,-1$.  They satisfy the same recurrence since
$p_{i+N}=p_i$.  Induction in both directions gives
\[
  X_{i+N}=X_i\quad\text{in }A/J_N\quad(i\in\Z).
\]
It follows that $[Z_{i\bmod N}]\mapsto[X_i]$ defines a homomorphism
$A_N\to A/J_N$.  It is inverse to the homomorphism induced by
$W_N$: one composite fixes all cyclic generators, and the other fixes
the classes of $x,y$.  This proves \eqref{eq:fixed-scheme-ideal}.

Proposition~\ref{prop:macro-shift} gives
$F^{n*}(x,y)=(X_N,X_{N-1})$.  The equalizer of $F^n$ and the identity
on the affine plane has exactly the ideal $J_N$, which proves the
scheme identification, not only an identification of its point set.
The intertwining formula follows on the generators $X_i$, and shifting
by $nk=N$ fixes every cyclic generator.  The dimension already proved
is thus the algebra length of this scheme, including multiplicities.
\end{proof}
